% Part: history
% Chapter: biographies 
% Section: alonzo-church

\documentclass[../../../include/open-logic-section]{subfiles}

\begin{document}

\olfileid{his}{bio}{chu}

\olsection{الونزو چرچ}

\olphoto{church-alonzo}{الونزو چرچ}

الونزو چرچ 14 جون 1903 نوں واشنگٹن ڈی سی وچ جمیا۔ نِکّی عمر وچ ہوا والی بندوق دے اک حادثے نال چرچ دی اک اکھ دی نظر چلی گئی۔ اوہنے 1920 وچ کنیکٹیکٹ وچ یونیورسٹی لئی تیاری والا سکول پورا کیتا تے اوسے سال پرنسٹن وچ یونیورسٹی دی پڑھائی شروع کیتی۔ اوہنے 1927 وچ اپنی ڈاکٹریٹ دی پڑھائی پوری کیتی۔ پردیس وچ دو کُ سال گزار کے چرچ پرنسٹن مُڑ آیا۔ چرچ بڑی شائستگی تے احتیاط لئی مشہور سی۔ تختے اُتے اوہدی لکھائی بے عیب ہُندی سی۔ ضروری کاغذاں نوں سنبھالن لئی اوہ احتیاط نال اوہناں اُتے ڈوکو سیمنٹ (اک شفاف گوند) دی تہہ لا دیندا سی۔ علمی کم توں ہٹ کے اوہ سائنس فکشن دے رسالے پڑھنا پسند کردا سی، تے جے لکھت وچ کوئی غلطی دِس پَندی تاں مدیران نوں لکھن توں جھجکدا نہیں سی۔

چرچ دیاں علمی کامیابیاں وڈیاں سن۔ اپنے شاگرداں سٹیفن کلینی تے بارکلے راسر نال رَل کے اوہنے مؤثر حساب دا نظریہ، لیمبڈا کیلکولس، بنایا۔ ایہ کم ایلن ٹیورنگ دی ٹیورنگ مشین دی بناوٹ توں آزاد سی۔ قابلِ حساب ہون دیاں دوویں تعریفاں برابر نیں، تے اوہناں توں اوہ گل نکلدی اے جیہنوں ہُن \emph{چرچ–ٹیورنگ دا دعویٰ} آکھدے نیں، یعنی قدرتی عدداں دا کوئی فنکشن مؤثر طور اُتے حساب جوگا اے اُدوں تے صرف اُدوں جدوں اوہ ٹیورنگ مشین (یا لیمبڈا کیلکولس) نال حساب جوگا اے۔ اوہنے اوہ گل وی ثابت کیتی جیہنوں ہُن \emph{چرچ دا قضیہ} آکھدے نیں۔ پہلے درجے دے فارمولیاں دے جائز ہون دا فیصلہ کرن والا مسئلہ حل نہیں ہو سکدا۔

چرچ وڈی عمر تک اپنا کم کردا رہیا۔ 1967 وچ اوہ پرنسٹن چھڈ کے یو سی ایل اے چلا گیا، جتھے اوہ 1990 وچ ریٹائر ہون تک پروفیسر رہیا۔ چرچ 1 اگست 1995 نوں 92 سال دی عمر وچ وفات پا گیا۔

\begin{reading}
چرچ دی مختصر سوانح عمری لئی \citet{EndertonND} ویکھو۔ لیمبڈا کیلکولس تے فیصلہ مسئلے، Entscheidungsproblem، (چرچ دے دعوے) بارے چرچ دیاں اصلی لکھتاں \citet{Church1936,Church1936a} نیں۔ \citet{Aspray1984} وچ 1930 دی دہائی دے پرنسٹن دے ریاضی داناں بارے چرچ نال اک انٹرویو درج اے۔ چرچ نے 1936 توں 1979 تک \emph{Journal of
Symbolic Logic} وچ کتاباں اُتے تبصریاں دی اک لڑی لکھی۔ ایہ سارے جان مکفارلین دی ویب سائٹ اُتے محفوظ نیں \citep{MacFarlane2015}۔
\end{reading}
\end{document}
